\def\thoigian{90}%--Thời gian
\de{ĐỀ THI HỌC KỲ II NĂM HỌC 2025-2026}{THPT Mạc Đĩnh Chi - Tp HCM}
\begin{center}
	\textbf{PHẦN 1 - TRẮC NGHIỆM}
\end{center}
\Opensolutionfile{ans}[ans/nlc]
%Câu 1
\begin{ex}%[0D8N1-2]%[Dự án đề KTTT Môn Toán HKII NH25-26 - Dot 9- Xuan Vy Pham]%[THPT Mạc Đĩnh Chi - Tp HCM]%[GVPB - Ngô Tất Thành]
	Một quán trà sữa phục vụ thực đơn gồm $4$ loại trà sữa khác nhau và $6$ loại topping khác nhau. Một khách hàng muốn chọn một combo gồm $1$ loại trà sữa và $1$ loại topping. Hỏi khách hàng đó có bao nhiều cách chọn một combo như vậy?
	\choice
	{$10$}	
	{\True $24$}	
	{$120$}	
	{$720$}	
	\loigiai{Chọn một trong $4$ loại trà sữa khác nhau có $4$ cách.\\
		Chọn một trong $6$ loại topping có $6$ cách.\\
		Vậy số cách chọn một combo gồm $1$ loại trà sữa và $1$ loại topping là $6 \cdot 4=24$ cách.
	}
\end{ex}
%Câu 2
\begin{ex}%[0H9H3-4]%[Dự án đề KTTT Môn Toán HKII NH25-26 - Dot 9- Xuan Vy Pham]%[THPT Mạc Đĩnh Chi - Tp HCM]%[GVPB - Ngô Tất Thành]
	Trong mặt phẳng tọa độ $Oxy$, góc giữa hai đường thẳng $d \colon \sqrt{3}x-y+2025=0$ và $\Delta \colon \sqrt{3}x+y+2026=0$ bằng
	\choice
	{\True $60^\circ$}	
	{$30^\circ$}	
	{$45^\circ$}	
	{$90^\circ$}
	\loigiai{Ta có $\cos \left(d,\Delta\right)=\dfrac{\left| \sqrt{3} \cdot\sqrt{3} -1 \cdot 1 \right|}{\sqrt{\left(\sqrt{3}\right)^2+(-1)^2} \cdot \sqrt{\left(\sqrt{3}\right)^2+1^2}}=\dfrac{1}{2}$ nên $\left(d,\Delta\right)=60^\circ$.}	
\end{ex}
%Câu 3
\begin{ex}%[0D7H3-1]%[Dự án đề KTTT Môn Toán HKII NH25-26 - Dot 9- Xuan Vy Pham]%[THPT Mạc Đĩnh Chi - Tp HCM]%[GVPB - Ngô Tất Thành]
	Tích các nghiệm thực của phương trình $\sqrt{2x^2+1}=\sqrt{x^2+5}$ bằng
	\choice
	{$0$}	
	{$-1$}	
	{$1$}	
	{\True $-4$}	
	\loigiai{Xét phương trình $\sqrt{2x^2+1}=\sqrt{x^2+5}$ (*).\\
		Bình phương hai vế của phương trình, ta được
		$2x^2+1=x^2+5$ hay $x^2-4=0$.\\ Suy ra $\hoac{&x=2\\&x=-2.}$\\
		Thay lần lượt hai giá trị trị $x=2$, $x=-2$ vào phương trình (*), ta thấy hai giá trị này đều thỏa mãn.\\
		Vậy phương trình (*) có hai nghiệm là $x=2$, $x=-2$ và tích hai nghiệm là $2 \cdot (-2) =-4$.}
\end{ex}
%Câu 4
\begin{ex}%[0H9H3-2]%[Dự án đề KTTT Môn Toán HKII NH25-26 - Dot 9- Xuan Vy Pham]%[THPT Mạc Đĩnh Chi - Tp HCM]%[GVPB - Ngô Tất Thành]
	Trong mặt phẳng toạ độ $Oxy$, cho tam giác $ABC$ với $A(1;3)$, $B(4;6)$, $C(0;4)$. Gọi $M$ là
	trung điểm của đoạn thẳng $BC$. Phương trình tổng quát của đường thắng $AM$ có dạng
	$ax-y+c=0$. Tính $P=a^2-c^2$.
	\choice
	{$P=4$}	
	{$P=1$}	
	{\True $P=3$}	
	{$P=8$}	
	\loigiai{$M$ là
		trung điểm của đoạn thẳng $BC$ nên 
		$$\heva{&x_M=\dfrac{x_B+x_C}{2}=\dfrac{4+0}{2}=2\\&y_M=\dfrac{y_B+y_C}{2}=\dfrac{6+4}{2}=5.}$$
		Vậy $M(2;5)$.\\
		Ta có $\overrightarrow{AM}=(1;2)$ nên đường thẳng $AM$ có một vectơ pháp tuyến là $\overrightarrow{n}=(2;-1)$.\\
		Phương trình tổng quát của đường thẳng $AM$ là
		$$2(x-1)-1(y-3) = 0 \Leftrightarrow 2x-y+1=0.$$
		Do đó $a=2$, $c=1$ nên $P=a^2-c^2=2^2-1^2=3$.}
\end{ex}
%Câu 5
\begin{ex}%[0D7N1-1]%[Dự án đề KTTT Môn Toán HKII NH25-26 - Dot 9- Xuan Vy Pham]%[THPT Mạc Đĩnh Chi - Tp HCM]%[GVPB - Ngô Tất Thành]
	Biểu thức nào sau đây là tam thức bậc hai?
	\choice
	{$f(x)=x+\dfrac{1}{x}$}	
	{$g(x)=2x-5$}	
	{\True $h(x)=x^2-2$}	
	{$t(x)=x^3+x-2$}	
	\loigiai{Biểu thức $h(x)=x^2-2$ là tam thức bậc hai.}
\end{ex}
%Câu 6
\begin{ex}%[0H9N5-1]%[Dự án đề KTTT Môn Toán HKII NH25-26 - Dot 9- Xuan Vy Pham]%[THPT Mạc Đĩnh Chi - Tp HCM]%[GVPB - Ngô Tất Thành]
	Trong mặt phẳng tọa độ $Oxy$, cho elip $(E) \colon \dfrac{x^2}{25}+\dfrac{y^2}{16}=1$. Độ dài trục lớn elip bằng
	\choice
	{$8$}	
	{\True $10$}	
	{$50$}	
	{$32$}	
	\loigiai{Elip $(E) \colon \dfrac{x^2}{25}+\dfrac{y^2}{16}=1$ có $a^2=25$ nên $a=5$.\\
		Độ dài trục lớn elip $A_1A_2=2a=2 \cdot 5=10$.}
\end{ex}
%Câu 7
\begin{ex}%[0D8H3-6]%[Dự án đề KTTT Môn Toán HKII NH25-26 - Dot 9- Xuan Vy Pham]%[THPT Mạc Đĩnh Chi - Tp HCM]%[GVPB - Ngô Tất Thành]
	Trong tiệm còn $5$ cái bánh khác nhau. Một khách hàng có bao nhiêu cách lựa chọn mua một số bánh trong số $5$ cái bánh đó (tính cả trường hợp khách hàng mua hết cả $5$ cái bánh và không mua cái bánh nào)?
	\choice
	{\True $32$}	
	{$5$}	
	{$21$}	
	{$10$}
	\loigiai{Số cách lựa chọn mua
		một số bánh trong số $5$ cái bánh (tính cả trường hợp khách hàng mua hết cả $5$ cái bánh và không mua cái bánh nào) là
		$$\mathrm{C}_5^0+\mathrm{C}_5^1+\mathrm{C}_5^2+\mathrm{C}_5^3+\mathrm{C}_5^4+\mathrm{C}_5^5=(1+1)^5=32.$$}	
\end{ex}
%Câu 8
\begin{ex}%[0D0H1-3]%[Dự án đề KTTT Môn Toán HKII NH25-26 - Dot 9- Xuan Vy Pham]%[THPT Mạc Đĩnh Chi - Tp HCM]%[GVPB - Ngô Tất Thành]
	Trong một hộp có $10$ viên bi khác nhau, gồm $6$ viên bi màu đỏ và $4$ viên bi màu xanh. Lấy ngẫu nhiên đồng thời $3$ viên bi từ hộp. Gọi A là biến cố: \lq \lq Lấy được ít nhất $2$ viên bi màu
	đỏ\rq \rq. Có bao nhiều kết quả thuận lợi cho biến cố $A$?
	\choice
	{$60$}	
	{$20$}	
	{$100$}	
	{\True $80$}	
	\loigiai{
		Chọn $3$ viên bi gồm $2$ bi đỏ, $1$ bi xanh có $\mathrm{C}_{6}^{2} \cdot \mathrm{C}_{4}^{1}$ cách.\\
		Chọn $3$ viên bi gồm $3$ bi đỏ có $\mathrm{C}_{6}^{3}$ cách.\\
		Vậy $n(A)=\mathrm{C}_{6}^{2} \cdot \mathrm{C}_{4}^{1}+\mathrm{C}_{6}^{3}=80$.
	}
\end{ex}
%Câu 9
\begin{ex}%[0H9N1-3]%[Dự án đề KTTT Môn Toán HKII NH25-26 - Dot 9- Xuan Vy Pham]%[THPT Mạc Đĩnh Chi - Tp HCM]%[GVPB - Ngô Tất Thành]
	Trong mặt phẳng toa độ $ Oxy$, cho hai điểm $A(1;3)$ và $B(5;9)$. Toa độ trung điểm $M$ của đoạn thẳng $AB$ là
	\choice
	{\True $(3;6)$}	
	{$(2;3)$}	
	{$(2;6)$}	
	{$(2;4)$}
	\loigiai{$M$ là
		trung điểm của đoạn thẳng $AB$ nên 
		$$\heva{&x_M=\dfrac{x_A+x_B}{2}=\dfrac{1+5}{2}=3\\&y_M=\dfrac{y_A+y_B}{2}=\dfrac{3+9}{2}=6.}$$
		Vậy $M(3;6)$.\\}	
\end{ex}
%Câu 10
\begin{ex}%[0D0N1-2]%[Dự án đề KTTT Môn Toán HKII NH25-26 - Dot 9- Xuan Vy Pham]%[THPT Mạc Đĩnh Chi - Tp HCM]%[GVPB - Ngô Tất Thành]
	Một hộp chứa $3$ quả cầu trắng được đánh số từ $1$ đến $3$. Lấy ngẫu nhiên đồng thời hai quả cầu từ hộp đó. Không gian mẫu $\Omega$ của phép thử là
	\choice
	{$\Omega =\left\{1;2;3\right\}$}	
	{$\Omega =\left\{\left(1;2\right);\left(2;1\right);\left(1;3\right);\left(3;1\right);\left(2;3\right);\left(3;2\right)\right\}$}	
	{\True $\Omega =\left\{\left(1;2\right);\left(1;3\right);\left(2;3\right)\right\}$}	
	{$\Omega =\left\{\left(1;1\right);\left(2;2\right);\left(3;3\right)\right\}$}
	\loigiai{Lấy ngẫu nhiên đồng thời hai trong $3$ quả cầu được đánh số từ $1$ đến $3$ nên không gian mẫu $\Omega$ của phép thử là $\Omega =\left\{\left(1;2\right);\left(1;3\right);\left(2;3\right)\right\}$.
	}	
\end{ex}
%Câu 11
\begin{ex}%[0D8N2-4]%[Dự án đề KTTT Môn Toán HKII NH25-26 - Dot 9- Xuan Vy Pham]%[THPT Mạc Đĩnh Chi - Tp HCM]%[GVPB - Ngô Tất Thành]
	Một tổ có $10$ học sinh. Hỏi có bao nhiêu cách chọn ra $2$ học sinh từ tổ đó để đi trực nhật?
	\choice
	{$5$}
	{$90$}
	{\True $45$}
	{$4$}
	\loigiai{Chọn ra $2$ học sinh trong $10$ học sinh từ tổ đó để đi trực nhật có $\mathrm{C}_{10}^2=45$ cách.}
\end{ex}
%Câu 12
\begin{ex}%[0H9H4-3]%[Dự án đề KTTT Môn Toán HKII NH25-26 - Dot 9- Xuan Vy Pham]%[THPT Mạc Đĩnh Chi - Tp HCM]%[GVPB - Ngô Tất Thành]
	Trong mặt phẳng tọa độ $Oxy$, cho đường tròn $(C) \colon (x-2)^2+(y-3)^2=25$. Tiếp tuyến của đường tròn $(C)$ tại điểm $A(6;0)$ đi qua điểm nào sau đây? 
	\choice
	{\True $M(3;-4)$}
	{$N(4;-3)$}
	{$P(4;3)$}
	{$Q(3;4)$}
	\loigiai{Đường tròn $(C) \colon (x-2)^2+(y-3)^2=25$ có tâm $I(2;3)$.\\
		Tiếp tuyến $d$ của đường tròn $(C)$ tại điểm $A(6;0)$ có phương trình là
		$$(2-6)(x-6)+(3-0)(y-0)=0 \Leftrightarrow -4(x-6)+3y=0 \Leftrightarrow 4x-3y-24=0.$$
		Vậy $d \colon 4x-3y-24=0$.\\
		Vì $4 \cdot 3 -3 \cdot (-4) -24=0$ nên $d$ đi qua $M(3;-4)$.}
\end{ex}

\Closesolutionfile{ans}
\begin{center}
	\textbf{PHẦN 2 - CÂU TRẮC NGHIỆM ĐÚNG SAI}
\end{center}
\Opensolutionfile{ans}[ans/ds]
\setcounter{ex}{0}
%%Câu 1
\begin{ex}%[0D0H2-4]%[Dự án đề KTTT Môn Toán HKII NH25-26 - Dot 9- Hoàng Tuấn]%[THPT Mạc Đĩnh Chi - Tp HCM]%[GVPB - Ngô Tất Thành]
	Lớp $10$A có $10$ học sinh đăng kí tham gia thi kéo co, gồm có $6$ học sinh nam và $4$ học sinh nữ. Giáo viên chọn ngẫu nhiên $5$ học sinh để lập một đội thi. Xét tính đúng sai của các khẳng định sau
	\choiceTF
	{\True Số các chọn $5$ học sinh từ $10$ học sinh là $252$}
	{Số cách chọn đội mà không có học sinh nữ là $720$}
	{\True Số cách chọn đội có đúng $2$ học sinh nữ là $120$}
	{Xác suất để đội được chọn có số học sinh nam nhiều hơn số học sinh nữ là $\dfrac{5}{7}$}
	\loigiai{
		\begin{itemchoice}
			\itemch Số cách chọn $5$ học sinh từ $10$ học sinh là tổ hợp chập $5$ của $10$.\\
			Do đó, ta có
			$$\mathrm{C}_{10}^5=\dfrac{10!}{5!\cdot 5!}=252.$$
			\itemch Số cách chọn đội trong đó không có học sinh nữ là tổ hợp chập $5$ của $6$.\\
			Do đó, ta có
			$$\mathrm{C}_{6}^5=\dfrac{6!}{5!\cdot 1!}=6.$$
			\itemch Số cách chọn đội để có đúng $2$ học sinh nữ là
			$$\mathrm{C}_{6}^3\cdot \mathrm{C}_{4}^2=20\cdot 6=120.$$
			\itemch Gọi $A$ là biến cố đội được chọn có số học sinh nam nhiều hơn số học sinh nữ.\\
			Ta xét các trường hợp sau
			\begin{itemize}
				\item Trường hợp 1: $3$ học sinh nam, $2$ học sinh nữ.\\
				Khi đó số cách chọn là 
				$$\mathrm{C}_6^3\cdot\mathrm{C}_4^2=20\cdot 6=120.$$
				\item Trường hợp 2: $4$ học sinh nam, $1$ học sinh nữ.\\
				Khi đó số cách chọn là 
				$$\mathrm{C}_6^4\cdot\mathrm{C}_4^1=15\cdot 4=60.$$
				\item Trường hợp 3: $5$ học sinh nam, $0$ học sinh nữ.\\
				Khi đó số cách chọn là 
				$$\mathrm{C}_6^5\cdot\mathrm{C}_4^0=6\cdot 1=6.$$
			\end{itemize}
			Suy ra tổng số cách chọn thoả mãn đề bài là 
			$$n\left(A\right)=120+60+6=186.$$
			Lại có $n\left(\Omega\right)=252$.\\
			Suy ra xác suất cần tìm là 
			$$\mathrm{P}(A)=\dfrac{186}{252}=\dfrac{31}{42}.$$
		\end{itemchoice}
	}
\end{ex}

%%Câu 2
\begin{ex}%[0H9H5-7]%[Dự án đề KTTT Môn Toán HKII NH25-26 - Dot 9- Hoàng Tuấn]%[THPT Mạc Đĩnh Chi - Tp HCM]%[GVPB - Ngô Tất Thành]	
	Trong mặt phẳng tọa độ $Oxy$, cho parabol $(P)$ có phương trình $y^2 = 8x$. Gọi $F$ là tiêu điểm và $\Delta$ là đường chuẩn của $(P)$. Gọi $d$ là đường thẳng đi qua $F$ và vuông góc với trục đối xứng của $(P)$. Đường thẳng $d$ cắt $(P)$ tại $2$ điểm phân biệt $A$ và $B$. Xét tính đúng sai của các khẳng định sau:
	\choiceTF
	{\True $(P)$ có tiêu điểm là $F\left(2;0\right)$}
	{Phương trình đường chuẩn của $(P)$ là $\Delta\colon x - 2 = 0$}
	{\True $AB = 8$}
	{Đường tròn $(C)$ có đường kính $AB$ tiếp xúc với đường chuẩn $\Delta$}
	\loigiai{
		\begin{itemchoice}
			\itemch Ta có $y^2=2px=8x\Rightarrow x=4$.\\
			Suy ra tiêu điểm của $(P)$ là $F\left(2;0\right)$.
			\itemch Phương trình đường chuẩn của $(P)$ là $x=-\dfrac{p}{2}=-2$, hay $x+2=0$.
			\itemch Đường thẳng $d$ đi qua $F\left(2;0\right)$ và vuông góc với trục $Ox$ nên có phương trình $x=2$.\\
			Thay vào $(P)$ ta có
			$$y^2 = 8 \cdot 2 = 16 \Rightarrow y = \pm 4.$$
			Khi đó $A(2; 4)$ và $B(2; -4)$, suy ra độ dài đoạn $AB$ là
			$$AB=\sqrt{\left(2-2\right)^2+\left(-4-4\right)^2}=8.$$
			\itemch Đường tròn $(C)$ đường kính $AB$ có tâm là $F\left(2;0\right)$ và bán kính $R = \dfrac{AB}{2} = 4$.\\
			Khoảng cách từ tâm $F\left(2;0\right)$ đến đường chuẩn $\Delta\colon x + 2 = 0$ là
			$$d\left(F, \Delta\right) = |2 + 2| = 4.$$
			Vì $d\left(F, \Delta\right) = R$ nên đường tròn tiếp xúc với đường chuẩn.
		\end{itemchoice}
		
	}
	
\end{ex}


\Closesolutionfile{ans}
\begin{center}
	\textbf{PHẦN 3 - CÂU TRẮC NGHIỆM TRẢ LỜI NGẮN}
\end{center}
\setcounter{ex}{0}
\Opensolutionfile{ans}[ans/kq]
%-----Câu 1
\begin{ex}%[0H9V3-8]%[Dự án đề KTTT Môn Toán HKII NH25-26 - Dot 9- Lien Nguyen]%[THPT Mạc Đĩnh Chi - Tp HCM]%[GVPB - Ngô Tất Thành]
	Trong mặt phẳng tọa độ $Oxy$ đơn vị trên các trục là kilômét, một trạm viễn thông được đặt tại vị trí $M\left(-\dfrac{15}{4};30\right)$. Một chiếc xe bus chuyển động thẳng đều với vectơ vận tốc $\overrightarrow{v}=(30;40)$. Biết tốc độ của xe bus chính là độ lớn của vectơ vận tốc $\overrightarrow{v}$. Lộ trình của xe bus đi qua điểm $K(10;15)$. Trạm viễn thông có phạm vi phủ sóng là hình tròn tâm $M$, bán kính $R$. Hành khách trên xe nhận được tín hiệu từ trạm trong khoảng thời gian liên tục $1$ giờ, sau đó mất tín hiệu. Hãy tính bán kính phủ sóng của trạm viễn thông này, làm tròn kết quả đến hàng đơn vị của kilômét.
	\shortans{32}
	\loigiai{
		\begin{center}
			\begin{tikzpicture}[scale=0.1, >=stealth, font=\footnotesize]
				% PHẦN 1: KHAI BÁO THAM SỐ VÀ TỌA ĐỘ
				% M(-15/4; 30), d: 4x - 3y + 5 = 0, R = sqrt(1025) approx 32,0156
				\def\mx{-3.75} 
				\def\my{30} 
				\def\radius{32.0156}
				\path
				(\mx,\my)coordinate (M)
				(12.25,18)coordinate (H) % Hình chiếu của M lên d
				(-2.75,-2)coordinate (A)% Giao điểm 1
				(27.25,38)coordinate (B) % Giao điểm 2
				(0,0) coordinate (O)
				(1,0) coordinate (O1)
				(0,1) coordinate (O2);
				\draw[->, line width=0.7pt] (-40,0) -- (45,0) node[below] {$x$};
				\draw[->, line width=0.7pt] (0,-15) -- (0,70) node[left] {$y$};
				\draw[orange!80!black, line width=0.6pt, dashed, name path=circle] (M) circle (\radius);
				\node[orange!80!black, font=\bfseries] at ($(M)+(55:35)$) {$(C)$};
				\draw[blue!60!black, line width=0.8pt] ($(A)!-0.4!(B)$) -- ($(A)!1.4!(B)$) node[right] {$d$};
				\draw[line width=0.5pt, dashed] (M) -- (A) node[midway, above left] {$R$};
				\draw[line width=0.5pt, dashed] (M) -- (B)(M)--(H);
				\pic[draw, angle radius=3mm,fill=pink]{right angle =B--H--M};
				\pic[draw, angle radius=3mm,fill=pink]{right angle =O1--O--O2};
				\foreach \p/\g in {M/150, H/-20, A/-80, B/-20,O/-55} {
					\draw [fill=black](\p) circle (2pt)node[shift={(\g:7pt)}]{$\p$};
				}
			\end{tikzpicture}
		\end{center}
		Ta có tốc độ của xe bus là $\left|\overrightarrow{v}\right|=\sqrt{30^2+40^2}=50$ (km/h).\\
		Quỹ đạo chuyển động của xe bus là đường thẳng $d$ đi qua $K(10;15)$ và có vectơ chỉ phương $\overrightarrow{u}=(30;40)$, hay $\overrightarrow{u}=(3;4)$.\\
		Do đó đường thẳng này có vectơ pháp tuyến $\overrightarrow{n}=(4;-3)$.\\
		Phương trình quỹ đạo xe bus là $4(x-10)-3(y-15)=0$, hay $4x-3y+5=0$.\\
		Gọi $H$ là hình chiếu vuông góc của $M$ lên đường thẳng $d$.\\
		Khoảng cách từ trạm $M\left(-\dfrac{15}{4};30\right)$ đến quỹ đạo xe bus là
		$$MH= \mathrm{d}=\dfrac{\left|4\cdot\left(-\dfrac{15}{4}\right)-3\cdot30+5\right|}{\sqrt{4^2+(-3)^2}}=\dfrac{100}{5}=20 \text{ (km)}.$$
		Trong $1$ giờ, xe đi được quãng đường $50$ km.\\
		Đây chính là độ dài dây cung $AB$ mà quỹ đạo xe cắt đường tròn phủ sóng. Nửa dây cung $AH=25$ km.\\
		Suy ra $R^2=20^2+25^2=1\,025$ (theo định lý Phythagore trong tam giác vuông $OAB$), nên $R=5\sqrt{41}\approx 32$.\\
		Vậy bán kính phủ sóng của trạm viễn thông là $32$ km.
	}
\end{ex}
%-----Câu 2
\begin{ex}%[0D0H2-7]%[Dự án đề KTTT Môn Toán HKII NH25-26 - Dot 9- Lien Nguyen]%[THPT Mạc Đĩnh Chi - Tp HCM]%[GVPB - Ngô Tất Thành]
	Một hộp có $8$ tấm thẻ được đánh số từ $1$ đến $8$. Lấy đồng thời $2$ tấm thẻ từ hộp đó. Tính xác suất để tích của hai số ghi trên $2$ tấm thẻ lấy ra là một số chẵn (\textit{làm tròn kết quả đến hàng phần trăm}).
	\shortans{0,79}
	\loigiai{
		Số cách lấy đồng thời $2$ tấm thẻ từ $8$ tấm thẻ là $\mathrm{C}_8^2=28$.\\
		Tích của hai số là số lẻ khi cả hai số đều là số lẻ. Trong các số từ $1$ đến $8$ có $4$ số lẻ là $1$, $3$, $5$, $7$.\\
		Số cách lấy được hai thẻ đều ghi số lẻ là $\mathrm{C}_4^2=6$.\\
		Vậy số cách lấy để tích hai số là số chẵn là $28-6=22$.\\
		Xác suất cần tìm là $P=\dfrac{22}{28}=\dfrac{11}{14}\approx 0{,}79$.\\
		Vậy xác suất cần tìm là $0{,}79$.
	}
\end{ex}
%-----Câu3
\begin{ex}%[0D7H1-2]%[Dự án đề KTTT Môn Toán HKII NH25-26 - Dot 9- Lien Nguyen]%[THPT Mạc Đĩnh Chi - Tp HCM]%[GVPB - Ngô Tất Thành]
	Có bao nhiêu giá trị nguyên của tham số $m$ để phương trình $(m-1)x^2+2(m-1)x+4=0$ vô nghiệm?
	\shortans{4}
	\loigiai{
		Xét phương trình $(m-1)x^2+2(m-1)x+4=0$.\\	
		Nếu $m=1$, phương trình trở thành $4=0$, nên phương trình vô nghiệm.\\
		Nếu $m\ne 1$, phương trình là phương trình bậc hai với $a=m-1$, $b=2(m-1)$, $c=4$.\\
		Khi đó biệt thức là $\Delta=\left[2(m-1)\right]^2-4(m-1)\cdot4=4(m-1)(m-5)$.\\
		Phương trình bậc hai vô nghiệm khi $\Delta<0$, tức là $4(m-1)(m-5)<0$.
		Suy ra $1<m<5$.\\
		Vì $m$ là số nguyên nên $m\in\{2;3;4\}$.\\
		Kết hợp với trường hợp $m=1$, ta có $m\in\{1;2;3;4\}$.\\
		Vậy có $4$ giá trị nguyên của $m$ để phương trình $(m-1)x^2+2(m-1)x+4=0$ vô nghiệm.
	}
\end{ex}
%-----Câu 4
\begin{ex}%[0D8H2-4]%[Dự án đề KTTT Môn Toán HKII NH25-26 - Dot 9- Lien Nguyen]%[THPT Mạc Đĩnh Chi - Tp HCM]%[GVPB - Ngô Tất Thành]
	Để thành lập một đội tình nguyện từ nhóm $10$ học sinh, giáo viên chọn ra $2$ học sinh để phân công làm Trưởng đoàn và Phó đoàn, sau đó chọn thêm $3$ học sinh khác từ số học sinh còn lại để làm thành viên. Hỏi giáo viên có tất cả bao nhiêu cách thực hiện?
	\shortans{5040}
	\loigiai{
		Chọn $2$ học sinh làm Trưởng đoàn và Phó đoàn là chọn có thứ tự, nên có $\mathrm{A}_{10}^{2}=10\cdot9=90$ cách.\\
		Sau khi chọn Trưởng đoàn và Phó đoàn, còn lại $8$ học sinh. Chọn thêm $3$ học sinh làm thành viên có $\mathrm{C}_8^3=56$ cách.\\
		Theo quy tắc nhân, số cách thực hiện là $90\cdot56=5\,040$.\\
		Vậy giáo viên có tất cả $5\,040$ cách thực hiện.
	}
\end{ex}
\Closesolutionfile{ans}
\begin{center}
	\textbf{PHẦN 4 - Phần tự luận}
\end{center}
\setcounter{ex}{0}
\Opensolutionfile{ans}[ans/tl]
%%-----Câu 1
\begin{ex}%[0D3H1-2]%[Dự án đề KTTT Môn Toán HKII NH25-26 - Dot 9- Lien Nguyen]%[THPT Mạc Đĩnh Chi - Tp HCM]%[GVPB - Ngô Tất Thành]
	Tìm tập xác định của hàm số $y=\dfrac{\sqrt{-x^2+2x}}{x-1}$.
	\loigiai{
		Hàm số xác định khi biểu thức dưới căn không âm và mẫu số khác $0$.\\
		Ta có điều kiện $-x^2+2x\ge 0$ và $x-1\ne 0$.\\
		Điều kiện $-x^2+2x\ge 0$ tương đương với $x(2-x)\ge 0$, suy ra $0\le x\le 2$.\\
		Mặt khác, $x-1\ne 0$ nên $x\ne 1$.\\
		Vậy tập xác định của hàm số là $\mathscr{D}=[0;1)\cup(1;2]$.
	}
\end{ex}
%%-----Câu 2
\begin{ex}%[0D8H3-4]%[Dự án đề KTTT Môn Toán HKII NH25-26 - Dot 9- Lien Nguyen]%[THPT Mạc Đĩnh Chi - Tp HCM]%[GVPB - Ngô Tất Thành]
	Tìm hệ số của số hạng chứa $x^3$ trong khai triển $(2x-1)^5$.
	\loigiai{
		Số hạng tổng quát trong khai triển $(2x-1)^5$ là $T_{k+1}=\mathrm{C}_5^k(2x)^k(-1)^{5-k}$ với $0\le k\le 5$.\\
		Số hạng chứa $x^3$ ứng với $k=3$.\\
		Khi đó hệ số của $x^3$ là $\mathrm{C}_5^3\cdot2^3\cdot(-1)^2=10\cdot8=80$.\\
		Vậy hệ số cần tìm là $80$.
	}
\end{ex}
%%-----Câu 3
\begin{ex}%[0H9V5-6]%[Dự án đề KTTT Môn Toán HKII NH25-26 - Dot 9- Lien Nguyen]%[THPT Mạc Đĩnh Chi - Tp HCM]%[GVPB - Ngô Tất Thành]
	Trong mặt phẳng tọa độ $Oxy$, cho tam giác $ABC$ với $A(2;6)$, $B(5;2)$ và $C(0;-3)$.
	\begin{enumerate}
		\item Viết phương trình tổng quát của đường trung trực của đoạn thẳng $BC$.
		\item Viết phương trình chính tắc của hypebol $(H)$ biết $(H)$ có độ dài trục thực bằng $2AB$, tiêu cự bằng $4BC$.
	\end{enumerate}
	\loigiai{
		\begin{enumerate}
			\item Ta có $B(5;2)$ và $C(0;-3)$ nên $\overrightarrow{BC}=(-5;-5)=-5\cdot (1;1)$.\\
			Trung điểm của đoạn thẳng $BC$ là $I\left(\dfrac{5}{2};-\dfrac{1}{2}\right)$.\\
			Vì đường trung trực của đoạn thẳng $BC$ vuông góc với $BC$, nên có vectơ pháp tuyến là $\overrightarrow{n}=(1;1)$.\\
			Phương trình đường trung trực của đoạn thẳng $BC$ là 
			Suy ra phương trình tổng quát là $1\cdot \left(x-\dfrac{5}{2}\right)+1\cdot\left(y+\dfrac{1}{2}\right)-2=0$, hay $x+y-2=0$.
			\item Ta có $AB=\sqrt{(5-2)^2+(2-6)^2}=\sqrt{3^2+(-4)^2}=5$.\\
			Lại có $BC=\sqrt{(5-0)^2+(2+3)^2}=\sqrt{25+25}=5\sqrt{2}$.\\
			Độ dài trục thực bằng $2AB$, nên $2a=2AB=10$, suy ra $a=5$.\\
			Tiêu cự bằng $4BC$, nên $2c=4BC=20\sqrt{2}$, suy ra $c=10\sqrt{2}$.\\
			Với hypebol chính tắc $\dfrac{x^2}{a^2}-\dfrac{y^2}{b^2}=1$, ta có $c^2=a^2+b^2$.\\
			Do đó $b^2=c^2-a^2=(10\sqrt{2})^2-5^2=200-25=175$.\\
			Vậy phương trình chính tắc của hypebol là $\dfrac{x^2}{25}-\dfrac{y^2}{175}=1$.
		\end{enumerate}
	}
\end{ex}

\Closesolutionfile{ans}